% !TeX root = ../main.tex
\chapter*{Abstract}
\addcontentsline{toc}{chapter}{Abstract}

% =========================================================================
% Draft Abstract preserved for integration upon adviser & panel approval:
% =========================================================================
% Under the constitutional doctrine of statutory preemption established in \textit{Magtajas v. Pryce Properties} and Section~5(a) of Republic Act No.~7160 (Local Government Code of 1991), local government ordinances in the Philippines are strictly subordinate to national statutes. An ordinance that prohibits an activity authorized by national law or authorizes an activity prohibited by national law is \textit{ultra vires} and void. In the Davao City Sangguniang Panlungsod, municipal legislative staff face substantial administrative and cognitive burdens when conducting manual ex-ante consistency reviews of proposed draft ordinances against decades of national laws and existing local measures. This challenge is intensified by the volume of positive law and the severe optical character recognition (OCR) degradation prevalent in historical municipal scanning archives.
% 
% To address these challenges, this study presents a coarse-to-fine semantic conflict detection architecture specifically designed for ex-ante legislative review and deployable within standard local government computing environments. The system decouples broad statutory candidate discovery from compute-intensive logical inference into a two-stage retrieve-then-entail pipeline. In Stage~1 (Coarse Candidate Retrieval), hybrid search combining sparse lexical matching (Okapi BM25) and dense neural embeddings (transformer bi-encoders) fused through Reciprocal Rank Fusion (RRF) with a statutory hierarchy priority safeguard searches a unified dual-corpus repository of 27,096 enactments (176,421 searchable provisions) to isolate a compact shortlist of top-$k$ candidate provisions ($k \le 50$). In Stage~2 (Fine Conflict Detection), candidate statute-ordinance pairs are classified through fine-tuned transformer cross-encoders into three directional normative categories: Contradiction, Entailment, or Neutral. High-risk contradictions are isolated using calibrated, cost-sensitive decision thresholds ($\tau^*$) and explained through intrinsic token-level self-attention heatmaps.
% 
% The computational framework was evaluated against an authoritative 350-pair ground truth benchmark curated with authentic Davao City legislative syntax, realistic local administrative bodies, and statutory drafting conventions, validated by an evaluation panel of fifteen active Sangguniang Panlungsod legal researchers ($\kappa \ge 0.61$). Empirical findings demonstrate that the hybrid Stage~1 retrieval pipeline elevates candidate recall and mean reciprocal rank substantially over unconstrained lexical baselines while executing in approximately 300~milliseconds on standard edge CPUs. In Stage~2, fine-tuned cross-encoders eliminate the affirmative entailment bias inherent in unadapted foundation models, achieving high contradiction detection sensitivity without requiring cloud-hosted supercomputers or generative large language models. The resulting system establishes an empirical baseline for low-resource Philippine legal informatics and provides an assistive decision-support tool for local legislative drafting.
% 
% \vspace{1em}
% \noindent\textbf{Keywords:} Semantic Conflict Detection, Ex-Ante Legislative Review, Statutory Preemption, Natural Language Inference, Information Retrieval, Cross-Encoder, Local Ordinances, Sangguniang Panlungsod.
